\chapter{The Identity Axis and Normalized Recurrence}
\label{chap:identity-axis}

\section{From a distinguished value to an axis}

The earlier chapters establish several exact appearances of the value
$\sigma=1/2$: complement fixedness, maximal binary entropy, equal-gain
cancellation, the Bloch equator, and the native PhaseNav closure condition.
The present chapter reorganizes those results around one geometric statement:
\emph{the half is an axis of identity rather than a frozen dynamical state}.

This terminology must be read precisely.  Let $X(\lambda)$ denote a state that
may evolve with an evolution parameter $\lambda$.  An identity axis is a
constraint or distinguished leaf that remains meaningful while coordinates
tangent to the leaf evolve.  In the binary reduction the transverse coordinate
is $\sigma$.  Define
\begin{equation}
 \Sigma_\star=\frac12,
 \qquad
 x=\sigma-\Sigma_\star.
 \label{eq:identity-centered-sigma}
\end{equation}
Then binary complement becomes reflection through the origin of the centered
chart:
\begin{equation}
 \sigma\mapsto1-\sigma
 \quad\Longleftrightarrow\quad
 x\mapsto-x.
 \label{eq:identity-reflection}
\end{equation}
The unique fixed point is $x=0$.  The system can nevertheless carry a
non-trivial phase, winding number, momentum or other tangent coordinate while
remaining on the balance leaf.

\begin{definition}[Binary identity axis]
\label{def:binary-identity-axis}
The binary identity axis is the fixed set of the complement involution in the
one-dimensional population coordinate,
\begin{equation}
 \Sigma_\star
 =\{\sigma\in[0,1]:\sigma=1-\sigma\}
 =\left\{\frac12\right\}.
\end{equation}
Its centered coordinate is $x=\sigma-1/2$.
\end{definition}

The definition is elementary.  The stronger interpretation that an analogous
$\Sigma$ class determines the persistent ``self'' of a general recurrent
system is a TIR/PhaseNav modelling principle, not a theorem about arbitrary
physical systems.  This monograph uses only the binary instance whose fixed
point is exact.

\section{Sign and zero as geometric relations}

Equation~\eqref{eq:identity-centered-sigma} gives a useful computational
consequence.  The sign of a centered scalar need not be introduced as an
independent semantic label.  It is the orientation of displacement relative to
the selected axis:
\begin{equation}
 \operatorname{sgn}_{\Sigma}(\sigma)
 =
 \begin{cases}
 -1,&\sigma<1/2,\\
 0,&\sigma=1/2,\\
 +1,&\sigma>1/2.
 \end{cases}
 \label{eq:identity-sign}
\end{equation}
Thus negative, zero and positive coordinates arise from one oriented chart.
Zero means vanishing displacement from the axis.  This statement is a fact
about the coordinate construction; it should not be inflated into the claim
that the real-number system itself has been physically derived.

The same idea extends to a phase torus.  If $q_1,q_2\in\mathbb R/\mathbb Z$,
the shortest oriented displacement may be represented in the interval
$[-1/2,1/2)$ by
\begin{equation}
 \Delta(q_1,q_2)
 =\left(q_2-q_1+\frac12\right)\bmod1-\frac12.
 \label{eq:shortest-turn-displacement}
\end{equation}
The signs of PhaseNav vector components can therefore be generated by geometry:
positive and negative values are opposite tangent orientations, and zero is an
identical phase coordinate.  This is the primitive used by the companion
CIEL/PhaseNav dependency package described in \cref{chap:geometry-dependencies}.

\section{The half-axis as the common zero of four residuals}

The preceding exact results can be restated as non-negative residuals.  For
$0<\sigma<1$, define
\begin{align}
 R_C(\sigma)
 &= (2\sigma-1)^2,\label{eq:residual-complement}\\
 R_H(\sigma)
 &= \ln2-\Hbin(\sigma),\label{eq:residual-entropy}\\
 R_I(\sigma)
 &= \left|\sqrt\sigma-\sqrt{1-\sigma}\right|^2,\label{eq:residual-interference}\\
 R_B(\sigma)
 &= \left|\exp[-2\pi i(1-\sigma)]+1\right|^2.
 \label{eq:residual-berry}
\end{align}
The first is the squared complement displacement.  The second is the entropy
deficit of \cref{chap:entropy}.  The third is the phase-$\pi$ cancellation
residual of \cref{chap:cancellation}.  The fourth measures deviation from the
$-1$ holonomy of one azimuthal loop in the qubit family of
\cref{chap:spinor}; Berry holonomy is standard geometric-phase structure
\citep{berry1984,simon1983}.

\begin{theorem}[Four-route half-axis cross-check]
\label{thm:four-route-half}
Within the domains and representations stated above,
\begin{equation}
 R_C(\sigma)=R_H(\sigma)=R_I(\sigma)=R_B(\sigma)=0
 \quad\Longleftrightarrow\quad
 \sigma=\frac12.
\end{equation}
Moreover, each residual separately has its relevant interior zero at
$\sigma=1/2$.
\end{theorem}

\begin{proof}
$R_C=0$ is equivalent to $2\sigma-1=0$.  Strict concavity of binary entropy
gives $R_H=0$ only at the unique maximum $1/2$.  The equal-gain cancellation
theorem gives $R_I=0$ only at equal populations.  Finally,
$R_B=0$ requires
$\exp[-2\pi i(1-\sigma)]=-1$; in the open unit interval the non-trivial
solution is $\sigma=1/2$.
\end{proof}

\begin{figure}[htbp]
 \centering
 \includegraphics[width=0.88\textwidth]{identity_half_routes.pdf}
 \caption{Four independently formulated non-negative residuals meet at the
 binary identity axis.  Their common zero is a consistency statement internal
 to the declared representations; it is not the missing zeta zero-state
 theorem.}
 \label{fig:identity-half-routes}
\end{figure}

The distinction in the caption is essential.  A common minimizer is valuable
because it reduces arbitrariness, but it cannot convert an unproved map from
zeta zeros to binary states into a theorem.

\section{Recurrence before angle}

A frequency can be defined without first introducing radians.  Let
\begin{equation}
 q\in\mathbb R/\mathbb Z
 \label{eq:normalized-cycle-q}
\end{equation}
be a normalized cycle coordinate.  A complete projective recurrence is
$q\mapsto q+1$, which is the same point in the quotient.  If $N(T)$ is the
signed winding count accumulated in elapsed parameter time $T$, define
\begin{equation}
 f_q
 =\lim_{T\to\infty}\frac{N(T)}{T}
 \label{eq:winding-frequency}
\end{equation}
when the limit exists.  For a finite interval the corresponding estimator is
$\Delta N/\Delta T$.

This definition is topological rather than angular.  Only after recurrence has
been identified do we choose an angular representation
\begin{equation}
 \phi=Cq,
 \label{eq:phase-closure-C}
\end{equation}
where $C>0$ is a closure period in the chosen angular unit.  Radian
parametrization is the conventional specialization
\begin{equation}
 C=2\pi.
 \label{eq:radian-closure}
\end{equation}
Then
\begin{equation}
 \dot\phi=C\dot q,
 \qquad
 \omega=2\pi f_q
\end{equation}
for uniform winding in standard radians.  The logical ordering is therefore
\begin{equation}
 \boxed{
 \text{recurrence}
 \longrightarrow
 \text{winding rate}
 \longrightarrow
 \text{angular representation}.
 }
 \label{eq:recurrence-order}
\end{equation}
This is useful for PhaseNav because a cycle counter can be implemented before a
particular angular convention is chosen.

\section{Projective and spinor return}

The double cover $\mathrm{SU}(2)\to\mathrm{SO}(3)$ supplies a second recurrence
scale.  If one projective loop is counted by $q\mapsto q+1$, the lifted spinor
changes sign after the first loop and returns after the second:
\begin{equation}
 +1\longrightarrow-1\longrightarrow+1.
 \label{eq:spinor-sheet-return}
\end{equation}
Consequently
\begin{equation}
 T_{\rm spin}=2T_{\rm proj}.
 \label{eq:spinor-projective-periods}
\end{equation}
No value of $\pi$ is needed to state the double-cover relation.  Under the
radian representation \eqref{eq:radian-closure}, the same structure is written
as the familiar $2\pi/4\pi$ return rule.

This separation also prevents a common category error.  The sign acquired by a
spinor under a physical $2\pi$ rotation is a representation-theoretic fact.
A Berry phase is a connection holonomy depending on a closed path.  They can
coincide at $-1$ for an appropriate equatorial loop, but the two notions are not
identical for arbitrary paths.

\section{Information cycle and the number twelve}

The TIR/Metatime layer uses the model assignment
\begin{equation}
 \frac{dI}{dq}=\frac{\ln2}{12}.
 \label{eq:info-per-projective-turn-soh}
\end{equation}
This equation is not a theorem of Shannon theory.  It says that the declared
information cycle distributes one binary information quantum over twelve
projective recurrences.  Once it is declared, however, the conversion to phase
is exact arithmetic.  Using \eqref{eq:phase-closure-C},
\begin{equation}
 \frac{dI}{d\phi}
 =\frac{\ln2}{12C}.
 \label{eq:info-angular-density-C}
\end{equation}
At standard radian closure,
\begin{equation}
 \left.\frac{dI}{d\phi}\right|_{C=2\pi}
 =\frac{\ln2}{24\pi}
 \equiv\kappa.
 \label{eq:kappa-cycle-crosswalk}
\end{equation}
Thus $\kappa$ factorizes into an arithmetic assignment and a geometric closure
scale.  This is a cross-derivation inside TIR, not an independent proof that the
number twelve is fundamental.  The independent derivation or falsification of
the twelve-cycle assignment remains an explicit debt.

\section{The 24/12/6 hierarchy}

The exact integer identity
\begin{equation}
 24=8\cdot3=12\cdot2=6\cdot4
 \label{eq:24-factorizations-soh}
\end{equation}
organizes three model decompositions used in the current programme.  The
arithmetic equality is exact.  The labels assigned to the factors are not:
\begin{align}
 8\cdot3&:\quad\text{eight declared mixing sectors times three flavours},\\
 12\cdot2&:\quad\text{twelve projective cycles times two half-turn units},\\
 6\cdot4&:\quad\text{six spinor cycles times four half-turn units}.
\end{align}
The second and third decompositions are mutually compatible with
\cref{eq:spinor-projective-periods}.  The first belongs to the wider TIR flavour
architecture and is cross-referenced, not proved, in this monograph.

\begin{figure}[htbp]
 \centering
 \includegraphics[width=0.82\textwidth]{cycle_hierarchy_24_12_6.pdf}
 \caption{Declared normalized-cycle hierarchy.  The counts are exact once the
 twelve-cycle model assignment is made; their physical interpretation remains
 model-level.}
 \label{fig:cycle-hierarchy}
\end{figure}

\section{Relation to the native theta bridge}

The native construction of \cref{chap:phasenav-native} already separates the
centered coordinate into
\begin{equation}
 z=\delta+it,
 \qquad
 \delta=\operatorname{Re}s-\frac12,
\end{equation}
where $\delta$ acts as antisymmetric radial gain shear and $t$ acts as phase
transport.  The present identity-axis formulation clarifies the same split:
\begin{itemize}
 \item the transverse displacement $\delta$ measures departure from the fixed
 identity axis;
 \item the tangent/phase coordinate can evolve while $\delta=0$;
 \item closure of the radial pair is therefore distinct from recurrence along
 the phase orbit.
\end{itemize}
This does not yet derive physical time from a phase period.  That proposed TIR
step is deliberately deferred and is not used in any theorem of this
monograph.

\section{Claim boundary}

The following items are exact or standard within their stated mathematical
settings: complement fixedness, the centered sign convention, normalized
recurrence, spinor double cover, the four-route half-axis cross-check, and the
algebraic conversion from a declared $dI/dq$ to $dI/d\phi$.

The following items remain model-level: interpreting $\Sigma$ as a universal
system-identity invariant, assigning one $\ln2$ cycle to twelve projective
recurrences, assigning the factors $8$ and $3$ to a complete physical mixing
architecture, and interpreting $\kappa$ as a fundamental physical coupling.

\status{The identity-axis and normalized-recurrence constructions sharpen the
internal geometry of the half.  They do not promote the open zero-state bridge
and do not prove RH.}
